%==================================================================
% Ini adalah bab 1
% Silahkan edit sesuai kebutuhan, baik menambah atau mengurangi \section, \subsection
%==================================================================

\chapter[PENDAHULUAN]{\\ PENDAHULUAN}

Minimal 3 halaman

\section{Latar Belakang}
\begin{sectioncontent}
    \hspace{\parindent}Uraian latar belakang penulisan laporan KKP (dalam bentuk paragraph), ceritakan latar belakang pemanfaatan teknologi IT pada tempat KKP
\end{sectioncontent}

\section{Identifikasi}
\begin{sectioncontent}
    Gambaran tentang apa yang ingin dicapai dalam pelaksanaan KKP
\end{sectioncontent}

\section{Batasan Masalah KKP}
\begin{sectioncontent}
    \begin{enumerate}[leftmargin=0.5cm]
        \item Jelaskan ruang lingkup dari KKP yang akan diangkat, berikan gambaran jaringan yang ada secara keseluruhan pada instansi KKP, kemudian sub jaringan atau layanan (service) yang dibahas secara detail (Tujuannya ialah untuk mendapatkan gambaran batas-batas penelitian)
        \item Batasan masalah. Jelaskan secara spesifik bagian sistem atau layanan yang akan dianalisa pada organisasi, dapat digunakan penguatan makna dengan mencantumkan sistem atau layanan yang tidak dibahas pada laporan
    \end{enumerate}
\end{sectioncontent}

\section{Sistematika Penulisan}
\begin{sistematika}
    \babsistematika{I}{PENDAHULUAN}{Berisi latar belakang masalah, identifikasi dan pembatasan masalah, perumusan masalah, tujuan dan manfaat penelitian, serta sistematika penulisan.}
    
    \babsistematika{II}{TINJAUAN ORGANISASI}{Menjelaskan tentang landasan teori yang digunakan dalam penelitian dan kerangka pemikiran.}
    
    \babsistematika{III}{PEMBAHASAN}{Menjelaskan metode penelitian yang digunakan, mulai dari pengumpulan data, analisis kebutuhan, perancangan sistem, hingga implementasi dan pengujian.}
    
    \babsistematika{IV}{PENUTUP}{Menyajikan hasil penelitian, serta pembahasan mengenai efektivitas dan manfaat penelitian yang dilakukan.}
\end{sistematika}